\subsection{Multiplication as Two-Layer Normalization: Convolution and Normal Forms}\label{subsec:emergent-arithmetic-stable-mul}

\begin{definition}[Stable multiplication]\label{def:stable-mul}
For $c,d\in X_\infty$, define
\[
c\otimes d:=\mathrm{Val}^{-1}\bigl(\mathrm{Val}(c)\cdot\mathrm{Val}(d)\bigr)\in X_\infty.
\]
\end{definition}

\begin{theorem}[Definitional isomorphism for multiplication (notation fixing)]\label{thm:mul-definitional}
$(X_\infty,\oplus,\otimes)$ is a commutative semiring, and $\mathrm{Val}$ yields a semiring isomorphism with $(\NN,+,\times)$.
\end{theorem}

\textbf{Proof:}
By Definitions~\ref{def:stable-add} and~\ref{def:stable-mul}, $\mathrm{Val}$ simultaneously preserves addition and multiplication. Since $\mathrm{Val}$ is a bijection, commutativity, associativity, and distributivity all transfer from $(\NN,+,\times)$, yielding the commutative semiring isomorphism.

Consequently, the natural numbers, addition, and multiplication are entirely defined within the stable address space by \textbf{syntactic constraints + normal form normalization}, without requiring $(+,\times)$ as pre-installed primitives.

\subsubsection{Multiplication arises internally from iterated addition}\label{subsec:mul-as-iterated-add}
\begin{theorem}[Multiplication arises from addition and iteration]\label{thm:mul-by-iterated-add}
Define on $X_\infty$
\[
c\odot d:=\underbrace{c\oplus c\oplus\cdots\oplus c}_{\mathrm{Val}(d)\text{ times}},
\qquad
\mathrm{Val}(d)=0\text{ gives }0_X,
\]
Then for any $c,d\in X_\infty$, $c\odot d=c\otimes d$.
\end{theorem}

\textbf{Proof:}
By Theorem~\ref{thm:add-definitional}, $\mathrm{Val}$ preserves $\oplus$ and is a bijection, so
\[
\mathrm{Val}(c\odot d)=\mathrm{Val}(c)\mathrm{Val}(d)=\mathrm{Val}(c\otimes d),
\]
whence $c\odot d=c\otimes d$. \qed

\begin{corollary}[First rigorous version of ``merging'']\label{cor:mul-no-new-primitive}
$\otimes$ is uniquely determined by $\oplus$ and iterable counting; it need not be introduced as a new primitive.
\end{corollary}

\subsubsection{Successor and folding: eliminating \texorpdfstring{$\oplus$}{oplus} as well}\label{subsec:succ-and-fold}
\begin{definition}[Successor and folding]\label{def:successor-fold}
Write $0_X:=(0,0,\dots)\in X_\infty$, and let $e_1\in\mathcal{D}$ be the first standard basis element. Define
\[
S:X_\infty\to X_\infty,\qquad S(c):=\Fold_\infty(c+e_1).
\]
\end{definition}

\begin{theorem}[Successor structure isomorphism]\label{thm:successor-structure}
$\mathrm{Val}$ yields a structure isomorphism:
\[
\mathrm{Val}(0_X)=0,\qquad \mathrm{Val}(S(c))=\mathrm{Val}(c)+1,
\]
and $\mathrm{Val}$ is a bijection.
\end{theorem}

\textbf{Proof:}
$\mathrm{Val}(0_X)=0$. Moreover,
\[
\mathrm{Val}(S(c))=\mathrm{Val}(c+e_1)=\mathrm{Val}(c)+\mathrm{Val}(e_1)=\mathrm{Val}(c)+1,
\]
since $\mathrm{Val}(e_1)=F_2=1$; bijectivity follows from Zeckendorf uniqueness. \qed

\begin{corollary}[Addition from iterated successor]\label{cor:add-from-successor}
Let
\[
c\oplus_S d:=S^{\circ\mathrm{Val}(d)}(c),
\]
then $c\oplus_S d=c\oplus d$, and $c\oplus d=\Fold_\infty(c+d)$.
\end{corollary}

\textbf{Proof:}
By Theorem~\ref{thm:successor-structure}, $\mathrm{Val}(c\oplus_S d)=\mathrm{Val}(c)+\mathrm{Val}(d)$; bijectivity gives $c\oplus_S d=c\oplus d$; the folding formula follows from Corollary~\ref{cor:add-as-fold}. \qed

\begin{theorem}[Multiplication arises from iterated successor]\label{thm:mul-from-successor}
Define
\[
c\otimes_S d:=\underbrace{c\oplus_S c\oplus_S\cdots\oplus_S c}_{\mathrm{Val}(d)\text{ times}},
\qquad \mathrm{Val}(d)=0\text{ gives }0_X,
\]
then $c\otimes_S d=c\otimes d$.
\end{theorem}

\textbf{Proof:}
By Corollary~\ref{cor:add-from-successor}, $\oplus_S=\oplus$, so
\[
\mathrm{Val}(c\otimes_S d)=\mathrm{Val}(c)\mathrm{Val}(d)=\mathrm{Val}(c\otimes d),
\]
and the conclusion follows by bijectivity. \qed

\begin{corollary}[Minimal primitive set]\label{cor:primitive-minimal-successor}
The arithmetic of the stable address layer is generated by the successor $S$ and the fold $\Fold_\infty$: $\oplus$ is the first-layer iterated readout of $S$, and $\otimes$ is the second-layer superposed readout of that iteration.
\end{corollary}

\subsubsection{Two-layer readout from the single primitive \texorpdfstring{$\circ$}{o}}\label{subsec:arith-composition}
This subsection provides a theorem-level version aligned with the chain ``scan iteration $\Rightarrow$ stable encoding $\Rightarrow$ normalization readout'': instead of pre-installing addition and multiplication as primitives, only \textbf{composition} is retained as the most fundamental primitive, and folding returns to the normal form in the stable address layer.
Take any nonempty set $A$ and a self-map $g:A\to A$. Write
\[
S_0:=\langle g\rangle=\{g^{\circ n}:n\in\NN\}\subseteq \mathrm{End}(A),\qquad
S_1:=\{I_n:n\in\NN\}\subseteq \mathrm{End}(\mathrm{End}(A)),
\]
where $I_n(f):=f^{\circ n}$. The sole operation on both layers is composition $\circ$.

\begin{theorem}[Two-layer readout from the single primitive \texorpdfstring{$\circ$}{o} simultaneously generates $(+,\times)$]\label{thm:composition-two-layer}
Define the embedding $\iota_0:\NN\to S_0$ by
\[
\iota_0(n):=g^{\circ n},\qquad g^{\circ 0}:=\mathrm{id}_A.
\]
and let $\iota_1:\NN\to S_1$ be $\iota_1(n):=I_n$. Then for any $m,n\in\NN$,
\[
\iota_0(m+n)=\iota_0(m)\circ \iota_0(n),
\qquad
\iota_1(mn)=\iota_1(m)\circ \iota_1(n).
\]
\end{theorem}

\textbf{Proof:}
By the definition of iteration and the associativity of composition,
\[
g^{\circ(m+n)}=g^{\circ m}\circ g^{\circ n},
\]
so $\iota_0(m+n)=\iota_0(m)\circ \iota_0(n)$. On the other hand, for any $f\in\mathrm{End}(A)$,
\[
(\iota_1(m)\circ \iota_1(n))(f)=I_m(I_n(f))=I_m(f^{\circ n})=(f^{\circ n})^{\circ m}
=f^{\circ (nm)}=I_{mn}(f).
\]
Since the equality holds for all $f$, we have $\iota_1(m)\circ \iota_1(n)=\iota_1(mn)$. \qed

\begin{theorem}[``De-primitivization'' implementation of the stable address semiring]\label{thm:arith-composition}
Using Definitions~\ref{def:stable-add} and~\ref{def:stable-mul}, and writing $\mathrm{Val}:X_\infty\to\NN$ for the value map, for any $(A,g)$ as in Theorem~\ref{thm:composition-two-layer}, define two ``process semantics'' maps
\[
\Phi_0:X_\infty\to \mathrm{End}(A),\quad \Phi_0(c):=g^{\circ\mathrm{Val}(c)},
\]
\[
\Phi_1:X_\infty\to \mathrm{End}(\mathrm{End}(A)),\quad \Phi_1(c):=I_{\mathrm{Val}(c)}.
\]
Then for any $c,d\in X_\infty$,
\[
\Phi_0(c\oplus d)=\Phi_0(c)\circ\Phi_0(d),\qquad
\Phi_1(c\otimes d)=\Phi_1(c)\circ\Phi_1(d).
\]
That is, $\oplus$ and $\otimes$ in the stable address layer are simultaneously realized by the single primitive $\circ$ acting on two layers of objects.
\end{theorem}

\textbf{Proof:}
By the definition of $\oplus$ and the first part of Theorem~\ref{thm:composition-two-layer},
\[
\Phi_0(c\oplus d)=g^{\circ\mathrm{Val}(c\oplus d)}
=g^{\circ(\mathrm{Val}(c)+\mathrm{Val}(d))}
 =g^{\circ\mathrm{Val}(c)}\circ g^{\circ\mathrm{Val}(d)}
=\Phi_0(c)\circ\Phi_0(d).
\]
Similarly, by the definition of $\otimes$ and the second part of Theorem~\ref{thm:composition-two-layer},
\[
\Phi_1(c\otimes d)=I_{\mathrm{Val}(c\otimes d)}
=I_{\mathrm{Val}(c)\mathrm{Val}(d)}
=I_{\mathrm{Val}(c)}\circ I_{\mathrm{Val}(d)}
=\Phi_1(c)\circ\Phi_1(d).
\]
\qed

\begin{corollary}[Composition-only readout falls back to stable addresses]\label{cor:composition-pullback}
Write $Z:=\mathrm{Val}^{-1}$, and define
\[
\rho_0:S_0\to X_\infty,\quad \rho_0(g^{\circ n}):=Z(n),\qquad
\rho_1:S_1\to X_\infty,\quad \rho_1(I_n):=Z(n).
\]
Then for any $c=Z(m),d=Z(n)\in X_\infty$,
\[
c\oplus d=\rho_0\bigl(\rho_0^{-1}(c)\circ \rho_0^{-1}(d)\bigr),
\qquad
c\otimes d=\rho_1\bigl(\rho_1^{-1}(c)\circ \rho_1^{-1}(d)\bigr).
\]
\end{corollary}

\textbf{Proof:}
By Theorem~\ref{thm:composition-two-layer}, $\rho_0(g^{\circ m}\circ g^{\circ n})=Z(m+n)$ and $\rho_1(I_m\circ I_n)=Z(mn)$; aligning with the definitions of $\oplus,\otimes$ yields the result. \qed

\begin{corollary}[Rigorous formulation of ``merging away $(+,\times)$'']\label{cor:primitive-free-arith}
At the configuration level, arithmetic primitives can be compressed to the following minimal generating set:
\begin{itemize}
  \item componentwise superposition of counting configurations (free commutative monoid structure);
  \item locally invertible rewriting induced by the Fibonacci relations (groupoid generators);
  \item selection of the normal form section $\Fold_\infty$ (or equivalently, $\mathrm{Val}$ and the Zeckendorf normal form $Z$).
\end{itemize}
At the process level, this can equivalently be expressed using the composition operation $\circ$ and the fold $\Fold_\infty$: $\oplus$ and $\otimes$ are respectively the readouts of composition on first- and second-layer objects, normalized via folding back to $X_\infty$.
Equivalently, the operation layer requires only $\circ$, and the visible layer requires only the encoding $Z$ (or $\mathrm{Val}/\Fold_\infty$).
\end{corollary}

\textbf{Proof:}
The three classes of primitives at the configuration level are directly given by Theorem~\ref{thm:monoid-quotient-is-N}, Theorem~\ref{thm:zeckendorf-transversal}, and Corollary~\ref{cor:fold-as-section}; the two-layer composition readout at the process level is given by Theorems~\ref{thm:composition-two-layer} and~\ref{thm:arith-composition}. The folding interface for addition is guaranteed by Corollary~\ref{cor:add-as-fold}, and the conclusion follows.

\begin{remark}[Quotient and remainder as normal form selection]\label{rem:division-as-selection}
Given $a,b\in X_\infty$ with $\mathrm{Val}(b)\ge 1$, the integer-level quotient-remainder relation
\[
\mathrm{Val}(a)=\mathrm{Val}(q)\,\mathrm{Val}(b)+\mathrm{Val}(r),\qquad 0\le \mathrm{Val}(r)<\mathrm{Val}(b)
\]
can be equivalently rewritten as an equation in the stable address layer:
\[
a=q\otimes b\ \oplus\ r,
\]
and the unique canonical representative $(q,r)$ is selected by the inequality constraint. Therefore ``division/quotient-remainder'' is not a new primitive but rather an \textbf{equation-solving + normal form selection} problem within the same rewriting groupoid; its computability is entirely inherited from the computability of $\Fold_\infty$.
\end{remark}

\begin{proposition}[Computable implementation]\label{prop:zeckendorf-arith-computable}
For any $c,d\in X_\infty$, $c\oplus d$ and $c\otimes d$ can be computed by a finite-step algorithm: first perform integer addition/multiplication at the value level to obtain $\mathrm{Val}(c)\pm\mathrm{Val}(d)$ or $\mathrm{Val}(c)\mathrm{Val}(d)$, then apply the Zeckendorf greedy algorithm to extract the unique normal form (Appendix~C).
\end{proposition}
